\section{Introducción}

El presente informe tiene por objetivo el diseño de los elementos estructurales estáticos del edificio en estudio: losas del piso tipo y vigas continuas ubicadas en el subterráneo. El análisis se enmarca en el proyecto de un edificio de hormigón armado de 21 pisos más un nivel de subterráneo.

En primer lugar, se realiza el predimensionamiento y diseño de armadura de las losas bidireccionales del piso tipo mediante el método de coeficientes, considerando las condiciones de apoyo de cada paño y diferenciando entre los niveles con hormigón G35 (piso 3) y G25 (pisos 4 en adelante). Se presentan los momentos de diseño, la compensación en bordes compartidos y los espaciamientos de armadura resultantes, junto con la verificación de deformaciones admisibles. En segundo lugar, se diseña la viga estática continua de cuatro tramos B4, B8, B15 y B16, ubicada en el subterráneo, determinando su armadura longitudinal a flexión y transversal a corte a partir de la envolvente de esfuerzos obtenida mediante modelación en SAP2000, y verificando las deformaciones de largo plazo conforme a ACI 318-19.
